% !TEX root = ../main.tex
% =====================================================================
% §5 Diagnostic Methodology & Honest Negative Findings — gdn2vessel benchmark-only D&B 稿
% 真源：STORY_FRAMEWORK.md（Claim 4 = 预登记诊断方法学 + 诚实负结果）
%        + PROJECT_LOG.md Entry 31/32（两命门死的完整证伪链）。
% 本节服务 Claim 4 / lever = D&B track 诚实评测贡献 + benchmark 判别力的正向证据。
% 三条负结果均已确证（实验在先、预登记 kill-criteria 写死）→ 照实写不 hedge。
% 但所有具体数字 \TODO 占位（禁臆造，verifier 核 csv 回填）；引用 \cite 占位。
% R6 严守：禁复活死命门当卖点（不写「delta 仍有微弱优势」「Frangi 门 DRIVE 涨说明有效」）。
% 双盲：无自引/机构；禁绝对化措辞。
% =====================================================================
\section{Diagnostic Methodology and Honest Negative Findings}
\label{sec:diagnostics}

A benchmark earns trust not only by ranking methods but by being able to
\emph{falsify} the hypotheses of those who build on it---including its own
authors. This section reports the diagnostic discipline that accompanies the
suite and the three negative findings it produced. We present these findings
as the intended output of the protocol rather than as failures to be hidden:
each is a case in which a method, a mechanism, or a metric was held to a
criterion fixed in advance and did not meet it, and each demonstrates that
the suite can detect exactly the kinds of limitation that a reconnection
evaluation is meant to surface.

\subsection{A pre-registered, falsify-the-crux diagnostic protocol}
\label{sec:diagnostics:protocol}

The diagnostic protocol that ships with this benchmark rests on two
principles. The first is \emph{pre-registration}: for every load-bearing
hypothesis, the kill-criteria---the numerical thresholds and the statistical
tests under which the hypothesis is to be rejected---are written down before
the corresponding data are seen. This guards against post-hoc rationalization
(hypothesizing after the results are known) and against the temptation to
relax a threshold once the numbers arrive. The second is a
\emph{three-layer line of defence} in which a hypothesis is first derived
semi-formally, then independently challenged by an adversarial review that
seeks to refute it, and finally checked numerically against the recorded
outputs, with no single layer able to confirm a claim on its own.

The discipline is deliberately ordered to test the load-bearing assumption
\emph{first}, with the cheapest decisive experiment, rather than last. A
recurring failure mode in method-driven work is to defer the crux---the one
assumption on which the headline rests---until after a large body of
supporting evidence has been assembled, at which point confirmation bias
makes an unfavourable result easy to read as favourable. By contrast, the
protocol here treats the crux as the thing to attempt to break at the
outset. The two negative findings reported below were each produced by a
pre-registered audit of this kind, applied to a hypothesis on which an
earlier framing of the project had rested. We report them in full because a
diagnostic methodology that only ever confirms its authors' hypotheses would
provide no evidence of discriminative power; the value of the protocol is
precisely that it rejected the hypotheses it was given. This follows the
spirit of benchmark work that documents negative and cautionary results
rather than suppressing them \cite{TODO_dandb_negative}.

\subsection{Negative finding I: associative memory is not mechanistically specific}
\label{sec:diagnostics:memory}

The first hypothesis subjected to a pre-registered audit was that a
delta-rule associative-memory module is mechanistically \emph{specific} for
reconnection---that its memory dynamics confer an advantage on tubular
reconnection over a plain stateful memory, and in particular that the memory
state can be made to bind same-root identity for re-identification. This
hypothesis was rejected on two independent grounds.

\paragraph{The mechanism is not specific (MQAR probe).}
We isolated the memory mechanism on a multi-query associative-recall
probe \cite{TODO_mqar}, a controlled retrieval task on which a
genuinely superior associative memory should separate from a baseline. Under
the pre-registered comparison, the delta-rule memory matched a plain gated
linear-attention memory at the only operating point at which both converged
(MQAR with sequence length $n{=}64$ and learning rate $10^{-3}$:
gdn2~$\approx$~\TODO{gdn2 acc, 核 csv} versus
gla~$\approx$~\TODO{gla acc, 核 csv}). At this point the two are
indistinguishable; the delta rule shows no advantage over a standard stateful
memory \cite{TODO_deltanet}. We emphasize what is \emph{not} concluded:
this is a negative result about \emph{specificity}, not a claim that either
memory is poor in absolute terms. The probe simply does not support the
hypothesis that the delta rule is the special ingredient.

\paragraph{The objective provides no gradient for the intended binding.}
The second ground is structural rather than empirical. In the reconnection
pipeline, the dominant training signal is the segmentation loss, and the
re-identification head is isolated from the memory state by three stop-gradient
boundaries. Consequently the segmentation objective exerts \emph{no} gradient
pressure to bind the memory state to a consistent same-root identity: the
property on which the ``memory performs re-identification internally''
hypothesis depended is never optimized for, and so cannot emerge from
training. Taken together, the non-specific probe result and the objective
mismatch establish that the memory module is one evaluated method on the
leaderboard, and nothing more. We make \emph{no} claim of mechanism
specificity, and the module's leaderboard standing is reported in
Section~\ref{sec:benchmark} on the same footing as every other baseline.

\subsection{Negative finding II: a differentiable vesselness gate does not robustly improve topology}
\label{sec:diagnostics:gate}

The second hypothesis was that a differentiable Frangi-style vesselness gate,
derived from the input image and used to modulate the memory, robustly
improves topological reconnection over an ungated variant. The gate adds a
negligible number of parameters relative to the backbone, so any improvement
could not be attributed to added capacity. The audit pre-registered a gate
that the gated variant had to pass to be accepted, with every threshold fixed
before the runs: (i) at least one dataset showing a clDice gap of at least
$+0.03$ in favour of the gated variant; (ii) a paired test with $p<0.05$ and
a bootstrap confidence interval whose lower bound lies above zero; and (iii)
a consistent direction of effect across both datasets.

The pre-registered gate was not passed. On DRIVE the gated variant was ahead
(clDice gap $\approx$~\TODO{DRIVE clDice gap, 核 csv}), but on CHASE the
direction reversed and the gated variant was behind
(clDice gap $\approx$~\TODO{CHASE clDice gap, 核 csv}). With the two datasets
pointing in opposite directions, and with neither dataset reaching the
pre-registered significance threshold
(\TODO{DRIVE p-value + bootstrap CI bound, 核 csv};
\TODO{CHASE p-value + bootstrap CI bound, 核 csv}), none of the three
pre-registered pass conditions was met. The honest reading is that the gate's
topological effect is \emph{not robust}: it is dataset-dependent rather than a
reliable improvement. In keeping with the pre-registered discipline, we did
not switch the evaluation axis, expand the seed count, or retune the gate's
scale to manufacture significance after the fact; a gate that passes only on
the dataset on which it happens to help is not a gate that improves topology.
The gated variant is therefore reported as one configuration the benchmark
evaluates, not as a contribution.

\subsection{Negative finding III: a success-rate metric can be saturated by over-prediction}
\label{sec:diagnostics:metric}

The third negative finding concerns a metric rather than a method, and is
detailed where the metric family is defined
(Section~\ref{sec:benchmark:srdiag}). In brief, the intuitive ``did the gap
close?'' success rate can be saturated by a degenerate over-predicting model,
because the underlying gap-closure test only checks whether any foreground
pixel covers the centre of the erased region; a model that floods the image
with foreground can be scored as closing every gap. We restate it here as the
third instance of the benchmark diagnosing a limitation---this time a
limitation of one of its own candidate metrics---which is why the primary
reconnection axis is a topological continuity measure and the success rate is
reported only in an over-prediction context.

\subsection{Negative findings as positive evidence of discriminative power}
\label{sec:diagnostics:positive}

The three findings above do not weaken the benchmark; they are the evidence
that it works. A reconnection evaluation suite is useful exactly to the extent
that it can show when a mechanism is not special, when a gating component is
unstable across datasets, and when a metric can be gamed. The suite did all
three: it found a memory mechanism to be non-specific against a plain
baseline, a vesselness gate to be unstable in the direction of its effect, and
a success-rate metric to be exploitable by over-prediction. None of these
would be visible without a fixed protocol, held-out evaluation, and
pre-registered criteria.

We are explicit about scope. These are established negative results, reported
as such; we do not soften them into qualified successes, and we do not
selectively report the dataset on which a rejected hypothesis happened to look
favourable. The modules involved---the associative-memory variant and the
gated variant---appear in the leaderboard as evaluated methods, on the same
footing as every other baseline, and not as headline contributions. Reporting
cautionary and negative findings of this kind is an intended use of a
benchmark, and dataset-and-benchmark venues explicitly invite it
\cite{TODO_dandb_negative}; surfacing where methods and metrics break is part
of what gives a benchmark its value.
